\documentclass[../../main.tex]{subfiles}
\graphicspath{{../../image/}{../../figure/}}
\begin{document}

\section{注记}\label{section:ch11-注记}

关于空间一维标量守恒律, 尤其真正非线性情形, 文献浩繁, 始于 Hopf [1] 关于 Burgers 方程的开创性论文, 前章已引。

1950 年代, 定性理论由以 Oleinik [1,2,4] 为首的俄罗斯学派发展, 基于消失黏性方法及 Lax-Friedrichs 有限差分格式(Lax [1])。正是在这一背景下,\cref{theorem:ch11-11-2-2} 最初建立。读者可在 Smoller [3] 的教材中找到阐述。该进路的顶点是多空间维标量守恒律理论的发展, 见第 VI 章。

在另一方向, Lax [2] 发现了 Cauchy 问题解的显式表示\eqref{align:ch11-11-4-10}, 并用它建立不变量的存在性(\cref{theorem:ch11-11-4-2})、紧支撑初值下 $N$ 波的发展(\cref{theorem:ch11-11-6-1})以及周期初值下锯齿剖面的形成(\cref{theorem:ch11-11-7-3})。于是齐次真正非线性标量守恒律的整套理论可建立在 Lax 公式之上。Oleinik [1] 推导了上述公式的推广, 适用于非齐次真正非线性标量守恒律。如\cref{chapter:ch11-一维标量守恒律}所述, 这些公式实际上是 Hamilton-Jacobi 方程粘性解的性质, 而非守恒律的性质, 因为它们可自然推广到前者的多空间维版本, 却不能推广到后者。


系统的、基于广义特征线的齐次标量守恒律理论由 Dafermos [4,12] 发展。Balian [1] 在离散框架中研究了熵不等式。Bressan [2,3,5,6,7,11], Bressan and Colombo [1], Bressan and Guerra [1], Bressan and Piccoli [1,2], Bressan and Russo [1], Colombo, Guerra and Mao [1], Colombo and Mauri [1], Guerra [1], Marsan [1] 以及 Risebro and Tveito [1] 给出了特殊的显式表示。

标量守恒律\eqref{align:ch11-11-1-1}的可容许解由子节 11.1--11.4 中的广义特征线刻画, 这是 Dafermos [8] 的贡献。Lax [3] 是引入并刻画分界线的第一人。
理论的实质性应用出现在: 研究 $L^1$ 收缩性与稳定性(Dafermos [7], [14]); 确定解的大时间行为(Dafermos [2,5,9,11], Dafermos and Hsiao [1], Hsiao [1]); 推导激波形成时间的显式公式(Dafermos [13]); 以及解的结构与正则性(Dafermos [6,10], Dafermos and Kecöeoghlu [1])。

$L^1$ 稳定性估计\eqref{align:ch11-11-8-12}来自 Dafermos [15]。\cref{chapter:ch14-前沿追踪与标准riemann半群}将其推广为方程组守恒律的 $L^1$ 稳定性。


通过广义特征线推导的激波形成显式公式出现在 Dafermos [13]。这一结果后来与 Enquist and Runborg [1] 以及 Liu [3,4] 的工作相联系。

在非线性输运方程(11.9.1)的框架下, 分界线、不变量与 Lax 公式的理论由 Dafermos [9,11] 与 Hsiao [1] 发展。

Oleinik [1] 是第一个研究齐次标量平衡律大时间行为的人。Lax [5] 引入并分析了特征线系统\eqref{align:ch11-11-10-2}与相关的时间周期驻波。周期零均值解渐近趋于有限振幅锯齿剖面的事实由 Dafermos [2] 发现。

非齐次守恒律的大时间行为由 Dafermos [5] 研究, 其通量 $f(u,x)$ 具有\eqref{align:ch11-11-11-1}中描述的结构, 临界点由极小与鞍点组成。

当通量 $f$ 具有拐点时, 可容许 $BV$ 解的结构在 Dafermos [10]、Liu [2] 与 Wagner [1] 中给出。$m$ 阶弱波与掠射线的出现由 Dafermos [10] 发现。分数阶有界变差 $u(\cdot,t)$ 的性质由 Chen, Dafermos, Wang and Wang [1] 给出。紧支撑初值下 $N$ 波渐近行为与周期零均值解的渐近估计来自 Dafermos [10]。

在空间一维标量守恒律情形, Bianchini and Marconi [1], 以及 De Lellis and Rivière [1] 精细化了\cref{chapter:ch06-标量守恒律L1理论}中概述的解结构结果, 证明 $L^\infty$ 解连续, 除可能在一维 Hausdorff 可求长集上, 熵产生集中于该集。11.13 节中连续解的讨论遵循 Dafermos [26]。应注意, 尽管仅连续的解在守恒律中罕遇, 但它们自然出现于一类带正则化源的平衡律设置中。有趣的代表是 Hunter-Saxton 方程与 Camassa-Holm 方程。本章所阐述技术精神下这些方程的处理见 Dafermos [31] 与 Bressan, Chen and Zhang [2]。关于非凸通量平衡律连续解的最新精细处理, 见 Alberti, Bianchini and Caravenna [1]。


空间一维标量守恒律与平衡律已知甚多, 欲全面覆盖殆不可能。下文仅是相关结果的抽样。

广义特征线的概率解释见 Rezakhaniou [1]。广义特征线方法在弹性静力学中, 不可压缩与不可伸长约束下的有趣应用见 Choksi [1]。进一步应用见 Dafermos [26] 与 Shearer and Dafermos [1]。

Lax 公式在四分之一平面上可容许解的显式表示, 与半平面对应, 由 LeFloch [1] 与 LeFloch and Nédélec [1] 给出。Lax 公式的类似物也已对激波曲线与稀疏波曲线重合的特殊系统推导; 见 Benzoni-Gavage [1]。

当 $g_\alpha(u)=f(u)v_\alpha$, 其中 $v$ 是常向量时, \eqref{align:ch11-11-2-1} 的类似物对多空间变量标量守恒律\eqref{align:ch06-6-1-1}成立(Hoff [1])。

输运距离中的收缩性质见 Esbelliorn, Gigli and Otto [1]。

标量守恒律的 Chapman-Enskog 型正则化见 Schochet and Tadmor [1]。另一引起相当注意的正则化是将标量守恒律与椭圆方程耦合, 产生所谓辐射气体方程组模型。见 Serre [20] 与 Lattanzio and Marcati [3]。

不同于\cref{chapter:ch06-标量守恒律L1理论}中讨论的一种动力学表示见 Brenier and Corrias [1]。另见 Vasseur [2]。


Panov [1] 以及独立地 De Lellis, Otto and Westdickenberg [2] 证明, 仅对一个一致凸熵的熵不等式即足以在 $L^\infty$ 中挑出唯一可容许弱解。

Besov 空间中解的正则性由 Lucier [2] 建立。数值格式收敛速率见 e.g. Nessyahu and Tadmor [1] 与 Osher and Tadmor [1]。

标量守恒律与"无压气体"系统\eqref{align:ch07-无压气体}及相关"黏性粒子"模型的联系由 E, Rykov and Sinai [1], Brenier and Grenier [1], Bouchut and James [1] 与 Tonon [1] 研究。无压气体的有趣理论由 Wang and Ding [1], Wang, Huang and Ding [1], Huang and Wang [1], Li and Warnecke [1], Ding and Huang [1], Huang [2] 与 Sever [9,10,11,12] 发展。广义特征线理论在诸多此类论文中起有用作用。

随机周期外力下的均匀化效应由 E [2,3], E and Serre [1], E, Khanin, Mazel and Sinai [1] 与 Amadori [2] 论证。

随机效应见 Holden and Risebro [1], Avellaneda and E [1], Bertoin [1], Jong Uhm Kim [1], Menon and Pego [1] 与 Chen, Ding and Karlsen [1]。

标量守恒律关联的边值问题与边界控制问题见 Andreianov and Sbihi [1,2] 与 Ancona and Marson [1,2]。

通量间断的情形由 Lyons [1], Risebro [2], Klingenberg and Risebro [1], Klausen and Risebro [1], Garavello, Natalini, Piccoli and Terracina [1], Chen, Even and Klingenberg [1], Diehl [1], Carrillo [1], Adimurthi and Gowda [1], Adimurthi, Dutta, Ghosal and Gowda [1], Adimurthi, Mishra and Gowda [1,2,3], Ostrov [2], Coclite and Risebro [1], Andreianov, Karlsen and Risebro [2], Holden, Karlsen and Mitrovic [1] 与 Crasta, De Cicco and de Philippis [1] 讨论。

如\cref{chapter:ch11-一维标量守恒律}所见, 当 $f$ 有拐点时, 解的结构因接触间断的形成而相当复杂, 接触间断成为向未来传播信号的源。解的一般结构由 Bianchini and Yu [2,3] 研究。另见 Dafermos [13], Jenssen and Sinestrari [1], Marson [1], Ballou [1,2], Guckenheimer [1], Chevry [4], Tang, Wang and Zhao [1], Jinghua Wang [1], Kim and Lee [1], Kim and Kim [1] 与 Ha and Kim [1]。大时间行为在 Dafermos [1,11], Greenberg and Tong [1], Conlon [1], Kuo Shung Cheng [1,2,3], Weinberger [1], Sinestrari [2], Yong Jung Kim [3,4] 与 Baiti and Jenssen [1] 中研究。特别地, 即使初值仅有界, $f'(u(\cdot,t))$ 在 $(-\infty,\infty)$ 上局部有界变差的证明见 Kuo Shung Cheng [3]。


在 $f(u)=u^m$ 特殊情形, 解的性质可借助底层自相似变换有效研究; 见 Bénilan and Crandall [1] 与 Liu and Pierre [1]。后者也考虑初值仅为测度的情形。该方向的发展见 Vasseur [6] 与 Chasseigne [1]。$m\to\infty$ 的极限行为见 Xiangsheng Xu [1]。

平衡律情形下源的正则化效应、解向行波的时间渐近收敛以及解的松弛极限, 由 Mascia [1,2,3] 与 Mascia and Terracina [1] 建立。另见 Isaacson and Temple [5]。

非齐次标量守恒律或平衡律可重构为两个守恒律或平衡律的共振系统处理; 见 Isaacson and Temple [5] 与 Amadori, Gosse and Guerra [2]。

其他相关发展见 Glass [2], Hayes and Shearer [2], Tadmor and Tang [1], Colombo and Goatin [1], Holden, Priuli and Risebro [1], Adimurthi, Dutta, Ghosal and Gowda [2,3,4], Andreianov and Seguin [1], Benzoni-Gavage [7], Coclite and Coclite [1,2], Ancona, Glass and Nguyen [1], Bourdarias, Gisclon and Junca [4], Colombo, Mercier and Rosini [1], Colombo and Rossi [1], Corli and Rohde [1], Bianchini and Modena [1] 与 Bank and Ben-Artzi [1]。

\end{document}
